% GENERATED FILE. Edit canonical lessons, then run npm run generate:books.
% canonical-source-hash: fnv1a64:b5e39ce9e074629e
% canonical-lessons: JA-C359-shakkin, JA-C359-rishi, JA-C359-tatekaeru, JA-C359-henkin, JA-C359-kingaku

\chapter{Borrowing and Lending}
\label{ch:ja-borrowing-and-lending}
\hlchaptermodality{359}

\begin{chapteropening}
\textbf{By the end of this chapter:} \emph{I can talk about owing money and paying it back.}

Complete the last lesson of chapter 359: I can talk about owing money and paying it back.
\end{chapteropening}

\section[\texorpdfstring{\ja{しゃっきん} (shakkin)}{しゃっきん (shakkin)}]{\ja{しゃっきん} (shakkin) — a debt}

\label{lesson:JA-C359-shakkin}



\begin{quote}
Before the new one: say the Japanese for to pay into an account, then the Japanese for a bank balance.
\end{quote}

\subsection*{You'll want to know: \ja{しゃっきん}}
\textbf{\ja{しゃっきん}} — \emph{shakkin} — \textquotedblleft{}a debt\textquotedblright{}.

\begin{quote}
\textbf{\ja{しゃっきんが} \ja{あったら、はやく} \ja{かえした} \ja{ほうが} \ja{いいです。}} — \emph{shakkin ga attara, hayaku kaeshita hō ga ii desu} — \textquotedblleft{}If you have a debt, you had better pay it back soon.\textquotedblright{}
\end{quote}

\begin{cousinweb}
\textbf{\ja{しゃっきん}} (\emph{shakkin}) is Sino-Japanese (kango): two Chinese characters read with their on'yomi, readings borrowed from Chinese. They are \emph{shaku}, \textquotedblleft{}borrow\textquotedblright{} (Middle Chinese \emph{tsjek}) and \emph{kin}, \textquotedblleft{}gold\textquotedblright{} (\emph{kim}) (KANJIDIC2; Middle Chinese in Baxter's transcription). The two join as \emph{shak-kin}.
\end{cousinweb}

\begin{grammarlens}[title={What you've built}]
One more word for saying what happens if.
\end{grammarlens}

\subsection*{Your turn}
\begin{itemize}
\raggedright
  \item \emph{Say it:} \emph{shakkin}
  \item \emph{Say it:} \emph{shakkin}, once more
  \item \emph{Say it:} \emph{shakkin ga attara, hayaku kaeshita hō ga ii desu} — If you have a debt, you had better pay it back soon.
  \item \emph{Recall:} say the Japanese for to pay into an account, then the Japanese for a bank balance, then say \emph{shakkin} again
\end{itemize}

\begin{tcolorbox}[breakable,colback=teal!4,colframe=teal!35!black,title={Before you move on}]
What does \ja{しゃっきん} mean? (\textquotedblleft{}A debt\textquotedblright{}.) Say it once more.
\end{tcolorbox}
\section[\texorpdfstring{\ja{りし} (rishi)}{りし (rishi)}]{\ja{りし} (rishi) — interest on money; interest on a loan}

\label{lesson:JA-C359-rishi}



\begin{quote}
Before the new one: say the Japanese for a bank balance, then the Japanese for a debt.
\end{quote}

\subsection*{You'll want to know: \ja{りし}}
\textbf{\ja{りし}} — \emph{rishi} — \textquotedblleft{}interest on money; interest on a loan\textquotedblright{}.

\begin{quote}
\textbf{\ja{おかねを} \ja{かりたら、りしを} \ja{はらいます。}} — \emph{okane o karitara, rishi o haraimasu} — \textquotedblleft{}If you borrow money, you pay interest.\textquotedblright{}
\end{quote}

\begin{cousinweb}
\textbf{\ja{りし}} (\emph{rishi}) is Sino-Japanese (kango): two Chinese characters read with their on'yomi, readings borrowed from Chinese. They are \emph{ri}, \textquotedblleft{}advantage\textquotedblright{} (Middle Chinese \emph{lijH}) and \emph{shi}, \textquotedblleft{}child\textquotedblright{} (\emph{tsiX}) (KANJIDIC2; Middle Chinese in Baxter's transcription).
\end{cousinweb}

\begin{grammarlens}[title={What you've built}]
One more word for saying what happens if.
\end{grammarlens}

\subsection*{Your turn}
\begin{itemize}
\raggedright
  \item \emph{Say it:} \emph{rishi}
  \item \emph{Say it:} \emph{rishi}, once more
  \item \emph{Say it:} \emph{okane o karitara, rishi o haraimasu} — If you borrow money, you pay interest.
  \item \emph{Recall:} say the Japanese for a bank balance, then the Japanese for a debt, then say \emph{rishi} again
\end{itemize}

\begin{tcolorbox}[breakable,colback=teal!4,colframe=teal!35!black,title={Before you move on}]
What does \ja{りし} mean? (\textquotedblleft{}Interest on money; interest on a loan\textquotedblright{}.) Say it once more.
\end{tcolorbox}
\section[\texorpdfstring{\ja{たてかえる} (tatekaeru)}{たてかえる (tatekaeru)}]{\ja{たてかえる} (tatekaeru) — to pay for someone (to be paid back)}

\label{lesson:JA-C359-tatekaeru}



\begin{quote}
Before the new one: say the Japanese for a debt, then the Japanese for interest on money.
\end{quote}

\subsection*{You'll want to know: \ja{たてかえる}}
\textbf{\ja{たてかえる}} — \emph{tatekaeru} — \textquotedblleft{}to pay for someone (to be paid back)\textquotedblright{}.

\begin{quote}
\textbf{\ja{さいふを} \ja{わすれたら、ともだちが} \ja{たてかえるでしょう。}} — \emph{saifu o wasuretara, tomodachi ga tatekaeru deshō} — \textquotedblleft{}If you forget your purse, a friend will probably pay for you.\textquotedblright{}
\end{quote}

Polite \textbf{\ja{たてかえます}}, \emph{tatekaemasu}; an ichidan verb. Past (polite): \textbf{\ja{あねが} \ja{りょうきんを} \ja{たてかえました。}}, \emph{ane ga ryōkin o tatekaemashita}.

\begin{cousinweb}
\textbf{\ja{たてかえる}} (\emph{tatekaeru}) is a native Japanese word (wago), inherited rather than borrowed from Chinese: it is built from \emph{tate}, \textquotedblleft{}setting up (tateru)\textquotedblright{} and \emph{kaeru}, \textquotedblleft{}to exchange\textquotedblright{}, and its kanji were attached later as kun'yomi (JMdict).
\end{cousinweb}

\begin{grammarlens}[title={What you've built}]
One more word for saying what happens if.
\end{grammarlens}

\subsection*{Your turn}
\begin{itemize}
\raggedright
  \item \emph{Say it:} \emph{tatekaeru}
  \item \emph{Say it:} \emph{tatekaeru}, once more
  \item \emph{Say it:} \emph{saifu o wasuretara, tomodachi ga tatekaeru deshō} — If you forget your purse, a friend will probably pay for you.
  \item \emph{Recall:} say the Japanese for a debt, then the Japanese for interest on money, then say \emph{tatekaeru} again
\end{itemize}

\begin{tcolorbox}[breakable,colback=teal!4,colframe=teal!35!black,title={Before you move on}]
What does \ja{たてかえる} mean? (\textquotedblleft{}To pay for someone (to be paid back)\textquotedblright{}.) Say it once more.
\end{tcolorbox}
\section[\texorpdfstring{\ja{へんきん} (henkin)}{へんきん (henkin)}]{\ja{へんきん} (henkin) — a refund}

\label{lesson:JA-C359-henkin}



\begin{quote}
Before the new one: say the Japanese for interest on money, then the Japanese for to pay for someone.
\end{quote}

\subsection*{You'll want to know: \ja{へんきん}}
\textbf{\ja{へんきん}} — \emph{henkin} — \textquotedblleft{}a refund\textquotedblright{}.

\begin{quote}
\textbf{\ja{しなものが} \ja{こわれたら、へんきんが} \ja{あるでしょう。}} — \emph{shinamono ga kowaretara, henkin ga aru deshō} — \textquotedblleft{}If the goods are broken, there will probably be a refund.\textquotedblright{}
\end{quote}

\begin{cousinweb}
\textbf{\ja{へんきん}} (\emph{henkin}) is Sino-Japanese (kango): two Chinese characters read with their on'yomi, readings borrowed from Chinese. They are \emph{hen}, \textquotedblleft{}return\textquotedblright{} (Middle Chinese \emph{pjonX}) and \emph{kin}, \textquotedblleft{}gold\textquotedblright{} (\emph{kim}) (KANJIDIC2; Middle Chinese in Baxter's transcription).
\end{cousinweb}

\begin{grammarlens}[title={What you've built}]
One more word for saying what happens if.
\end{grammarlens}

\subsection*{Your turn}
\begin{itemize}
\raggedright
  \item \emph{Say it:} \emph{henkin}
  \item \emph{Say it:} \emph{henkin}, once more
  \item \emph{Say it:} \emph{shinamono ga kowaretara, henkin ga aru deshō} — If the goods are broken, there will probably be a refund.
  \item \emph{Recall:} say the Japanese for interest on money, then the Japanese for to pay for someone, then say \emph{henkin} again
\end{itemize}

\begin{tcolorbox}[breakable,colback=teal!4,colframe=teal!35!black,title={Before you move on}]
What does \ja{へんきん} mean? (\textquotedblleft{}A refund\textquotedblright{}.) Say it once more.
\end{tcolorbox}
\section[\texorpdfstring{\ja{きんがく} (kingaku)}{きんがく (kingaku)}]{\ja{きんがく} (kingaku) — a sum of money; an amount}

\label{lesson:JA-C359-kingaku}



\begin{quote}
Before the new one: say the Japanese for to pay for someone, then the Japanese for a refund.
\end{quote}

\subsection*{You'll want to know: \ja{きんがく}}
\textbf{\ja{きんがく}} — \emph{kingaku} — \textquotedblleft{}a sum of money; an amount\textquotedblright{}.

\begin{quote}
\textbf{\ja{きんがくが} \ja{かわったら、れんらくします。}} — \emph{kingaku ga kawattara, renraku shimasu} — \textquotedblleft{}If the amount changes, I will let you know.\textquotedblright{}
\end{quote}

\begin{cousinweb}
\textbf{\ja{きんがく}} (\emph{kingaku}) is Sino-Japanese (kango): two Chinese characters read with their on'yomi, readings borrowed from Chinese. They are \emph{kin}, \textquotedblleft{}gold\textquotedblright{} (Middle Chinese \emph{kim}) and \emph{gaku}, \textquotedblleft{}sum\textquotedblright{} (\emph{ngaek}) (KANJIDIC2; Middle Chinese in Baxter's transcription).
\end{cousinweb}

\begin{grammarlens}[title={What you've built}]
One more word for saying what happens if.
\end{grammarlens}

\subsection*{Your turn}
\begin{itemize}
\raggedright
  \item \emph{Say it:} \emph{kingaku}
  \item \emph{Say it:} \emph{kingaku}, once more
  \item \emph{Say it:} \emph{kingaku ga kawattara, renraku shimasu} — If the amount changes, I will let you know.
  \item \emph{Recall:} say the Japanese for to pay for someone, then the Japanese for a refund, then say \emph{kingaku} again
\end{itemize}

\begin{tcolorbox}[breakable,colback=teal!4,colframe=teal!35!black,title={Before you move on}]
What does \ja{きんがく} mean? (\textquotedblleft{}A sum of money; an amount\textquotedblright{}.) Say it once more.
\end{tcolorbox}
